\section[What the subsequent fits identify]{What the subsequent exact-sample fits identify}
\label{sec:neutra-october-fit-attribution}

The subsequent experiment on randomized two- and three-component Gaussian
mixtures isolates a failure that occurs even before an approximate sampler
supplies the training data. Its four independent exact-sample banks contain
$4{,}096$ observations in total. The forward fit resamples these observations
with their empirical probabilities, so its expected minibatch gradient is
the gradient of the empirical forward cross entropy. Exact sampling removes
an approximate sampler as the cause of this failure; it does not remove
finite-data error or establish that the optimizer reaches the population
optimum.

In the October 4 experiment, both widths, $32$ and $64$, retained three IAF
stages. They received $4{,}096$ and $8{,}192$ forward updates respectively,
followed by $256$ reverse-KL updates with a fresh Adam state. The objective
switch occurred regardless of the forward assessment. Changing width and
update count together prevents this comparison from identifying a width
effect. Moreover, the learning histories did not establish a stationary
forward optimum. All four final maps failed the declared distribution screen,
and the native teacher experiments consequently did not run.

Checks of the eight saved forward and final maps found inverse and
log-determinant discrepancies of approximately $10^{-14}$ and directional
parameter-derivative discrepancies up to $1.29\times10^{-7}$ in the
checked FP64 reference calculations. No gradient clipping occurred during
these fits, and the sampled conditional-scale diagnostics did not show strong
cap saturation. These observations weaken specific explanations based on
incorrect local derivatives or clipping. They do not prove every execution
path correct, establish sufficient capacity, or validate the complete
training procedure.

The forward maps already assigned excessive mass away from the components.
For a two-dimensional Gaussian mixture
$p(x)=\sum_{k=1}^K w_k\varphi_k(x)$, let
$d_k^2(x)=(x-\mu_k)^\mathsf T\Sigma_k^{-1}(x-\mu_k)$ and define
\begin{equation}
 E=\{x:d_k^2(x)>8\ \text{for every }k\}.
\end{equation}
Conditional on component $k$, this event is contained in
$\{d_k^2(X)>8\}$. Since a component's squared Mahalanobis distance is
$\chi^2_2$, the exact target therefore satisfies
\begin{equation}
 p(E)=\sum_k w_k\Pr_{\varphi_k}(E)
 \leq\sum_k w_k e^{-4}=e^{-4}\simeq0.01832.
 \label{eq:neutra-mixture-ellipse-bound}
\end{equation}
The width-$64$ forward maps assigned approximately $8.39$ percent and
$21.62$ percent to this event for the two- and three-component targets.
These estimates use $32{,}768$ new Gaussian-base draws per target. Their
uncertainty concerns numerical integration for fixed fitted maps, rather
than variability across repeated training runs. The bound is an upper
bound on target mass, not a formula for its exact value.

The subsequent coverage loss has an exact reverse-KL interpretation that
does not require disjoint Gaussian components. Define the target
responsibilities and the learned responsibility masses by
\begin{equation}
 r_k(x)=\frac{w_k\varphi_k(x)}{p(x)},\qquad
 a_k=\int q(x)r_k(x)\,dx,\qquad
 q_k(x)=\frac{q(x)r_k(x)}{a_k}.
\end{equation}
Here $q_k$ is the conditional density induced by assigning target
responsibilities to draws from $q$; it is not assumed to be Gaussian.
The augmented densities
$\widetilde p(x,k)=p(x)r_k(x)=w_k\varphi_k(x)$ and
$\widetilde q(x,k)=q(x)r_k(x)=a_kq_k(x)$ have ratio $q(x)/p(x)$.
Summing the augmented KL over $k$ therefore leaves the marginal KL
unchanged and yields
\begin{align}
 D_{\mathrm{KL}}(q\|p)
 &=\sum_k\int a_kq_k(x)
       \log\frac{a_kq_k(x)}{w_k\varphi_k(x)}\,dx\\
 &=\sum_k a_k\log\frac{a_k}{w_k}
   +\sum_k a_kD_{\mathrm{KL}}(q_k\|\varphi_k).
 \label{eq:neutra-responsibility-kl-decomposition}
\end{align}
The first term penalizes incorrect component allocation. The second
penalizes conditional shape error, weighted by the probability that the
learned density itself assigns to that component.

\begin{table}[htbp]
 \centering
 \begin{tabular}{lrr}
  \toprule
  Quantity, width-$64$ three-component fit & Forward endpoint & After RKL\\
  \midrule
  Second responsibility mass $a_2$ & $0.21572$ & $0.02240$\\
  Allocation KL & $0.00319$ & $0.20984$\\
  Weighted conditional-shape KL & $1.69401$ & $0.41727$\\
  Total reverse KL & $1.69721$ & $0.62711$\\
  Heldout forward KL & $0.42979$ & $1.14216$\\
  \bottomrule
 \end{tabular}
 \caption{The objective switch reduces reverse KL while worsening coverage.
 The target second-component weight is $0.249591$. Estimates describe one
 pair of fixed checkpoints; they do not rank training methods.}
 \label{tab:neutra-objective-switch-attribution}
\end{table}

Table~\ref{tab:neutra-objective-switch-attribution} shows that the reduction
in weighted shape cost exceeds the increase in allocation cost. With common
integration draws, the change in $a_2$ is $-0.19333$, with an approximate
$95$ percent integration interval $[-0.19753,-0.18913]$. The reverse-KL
change is $-1.07010$, with estimated integration standard error $0.03867$.
Thus this coverage regression is compatible with correctly computed
reverse-KL optimization. It does not show that reverse KL always loses modes.
The poor forward endpoints also show why simply omitting the reverse stage
would leave a substantial fitting problem.

\section[Published remedies]{Remedies in the literature}
\label{sec:neutra-literature-remedies}

The combined objective in \eqref{eq:controlled-joint-loss} has a direct
precedent in \citet[][Methods, Eq.~(9)]{noe2019boltzmanngenerators}.
Their original double-well notebook performs maximum-likelihood pretraining
and then retains both maximum-likelihood and energy terms. Similarly,
\citet[][Sec.~III.B and Appendix~A.2]{gabrie2022adaptiveflows} describe
combining the two objectives after a sufficiently accurate forward fit.
Equation~\eqref{eq:controlled-kl-partition} explains the protection that
the forward term supplies: eliminating a region with positive target mass
has an unbounded forward penalty. Finite weights, finite examples and
nonconvex optimization still permit substantial error. In particular, the
earlier eight-case experiment in this chapter already tried joint and
joint-continuation training without attaining fresh posterior confirmation.
The literature supports testing this mechanism carefully; it does not make
it a previously untried or sufficient cure for the present failure.

Representation is a separate choice. A scalar affine transformation of a
Gaussian remains Gaussian at fixed conditioning variables. The nonlinear
monotone transformations of \citet[][Sec.~3 and Fig.~5]{huang2018naf}
can instead create multiple scalar density peaks through changes in their
derivative. Composing and permuting affine autoregressive stages can also
produce complex multivariate densities, so this distinction is not a proof
that the current IAF family cannot represent the mixture. It identifies a
reason to compare representations after controlling optimization effort.
The Gaussian mixtures considered here are positive throughout
$\mathbb R^d$; an impossibility argument based on disconnected support does
not apply to them. The exact transport in
\eqref{eq:controlled-exact-transports} is also a constructive reminder that
Gaussianization and numerical ease of learning that Gaussianization are
different questions.

\citet[][Sec.~3.1, Eqs.~(4)--(8)]{durkan2019splines} offer a useful
computational alternative. Their monotone rational-quadratic scalar splines
have analytic derivatives and inverses. In an autoregressive construction
the Jacobian remains triangular, with
\begin{equation}
 \log|\det J_T(z)|
 =\sum_j\log\left|\frac{\partial T_j(z)}{\partial z_j}\right|.
\end{equation}
The inverse within a spline bin follows from a quadratic equation, as in
the original implementation. This avoids a dense determinant, although an
autoregressive inverse can still require sequential coordinate evaluation.
The usual splines are continuously differentiable; regularity of the
transformed force at knots requires separate attention for HMC. Such a
comparison would evaluate an explicitly different map family, rather than
silently redefining the canonical IAF.

\citet[][Sec.~IV.E]{gabrie2022adaptiveflows} also recommend informed base
distributions, appropriate coordinate scaling and mixtures of maps fitted
to different modes. Their Gaussian-mixture experiment uses six pairs of
RealNVP coupling layers with width-$100$ conditioners. It is therefore not
an experiment with our three-stage IAF. Appendix~G.1 reports a residual
connection between modes even in the successful example and states that
additional training would make it thinner. The proposed mixture-of-maps
strategy reduces the burden on any individual map, but produces a mixture
proposal, not automatically a single invertible Gaussian coordinate chart
for the existing NeuTra HMC procedure.

Fresh training information addresses another failure mechanism. Gabri\'e's
Algorithm~1 alternates local moves, global flow proposals and forward
training. Local moves explore details the flow misses, while accepted
global proposals help adjust relative mode weights. Section~IV.C assumes
initial representatives of the relevant modes; the method does not promise
to discover every unrepresented basin. Moreover, the paper's ULA experiments
retain discretization error, even though corrected MALA is available in the
algorithm. Fixed-kernel invariance and convergence of a finite adaptive
implementation must be assessed separately.

Annealed flow transport divides a difficult transformation into transitions
between intermediate distributions. The practical AFT procedure of
\citet[][Sec.~3.3 and Appendix~F]{arbel2021aft} uses training, validation
and test populations to limit overfitting and separate fitting from
evaluation. \citet[][Sec.~2.3]{matthews2022craft} identify the remaining
sample-replenishment problem: repeatedly optimizing on one finite
population can exhaust the information available for learning a flow.
CRAFT uses repeated annealed passes with fresh initial particles and retains
importance, resampling and Markov corrections. Its finite-particle gradient
is not generally an unbiased gradient of the exact-target objective;
Section~3.2 instead gives an unbiased-gradient interpretation for a modified
objective involving the expected finite-particle distribution. These methods
improve the sampling-and-learning procedure, but cannot alone explain or
repair failure of a separate student already supplied with exact samples.

FAB strengthens the penalty for missed mass through an objective
proportional to $\int p(x)^2/q(x)\,dx$, with AIS targeting its normalized
integrand \citep[][Sec.~3.1]{midgley2023fab}. The method concentrates effort
where the target density is appreciable but the flow density is small.
Its applicability requires a finite integral and effective intermediate
transitions. The favorable dimensional-scaling analysis in that section
assumes factorized distributions and perfect intermediate transitions.
It does not resolve the previously observed q20 tail-applicability and
transition-calibration problems for an arbitrary starting map.

For rare transition regions, deliberate sampling is often necessary.
The reaction-coordinate objective of \citet{noe2019boltzmanngenerators}
encourages configurations that ordinary equilibrium sampling seldom visits.
Umbrella sampling instead uses overlapping biased windows; the EMUS method
of \citet[][Secs.~II--III, Eqs.~(9)--(17)]{thiede2016emus} estimates their
relative normalizations from overlap and combines their weighted information.
A useful reaction coordinate, sufficient overlap and adequate sampling within
each window remain assumptions. Correct reweighting improves information
about a rare region without assigning it an artificial posterior mass.
As the partition calculation above shows, this still leaves a rare region
with a small contribution to ordinary KL. Accurate derivatives there do not
follow from a small integrated density loss.

\section[Reproduction and transfer]{Reproducing training before testing transfer}
\label{sec:neutra-original-code-comparison}

The available original-code evidence has different scopes. The FAB audit
executed pinned author JAX operations, compatibility-adjusted upstream tests
and four matched optimizer iterations in each of four sampling/replay
configurations. The FP64 stress comparison passed $1{,}804$ checks. The final
compiled-callback TF32 comparison failed seven declared gradient checks.
These are meaningful finite numerical comparisons; they do not reproduce
converged training on the current separated mixtures. The precision-screen
failure also does not explain an FP64 exact-teacher run in which FAB never
executed.

The Gabri\'e reference executed selected original MH/MALA transitions and
target-density checks, but did not train the original RealNVP through its
complete adaptive controller. The available AFT/CRAFT work inspected author
source and checked local gradients and update ordering, without reproducing
complete original-controller training on these targets. NeuTra architecture
and derivative tests likewise do not reproduce an entire published training
experiment. Consequently, failure of the common IAF student cannot be
counted as failure of each of these literature methods.

A discriminating comparison begins with an author's documented target,
architecture, initialization, objective and sampling cadence. A matched
port then preserves those choices while comparing densities, gradients,
optimizer state and full training behavior. Substituting the canonical IAF
is a subsequent scientific change; switching objectives is another.
If the original and its matched port agree but fail after the substitution,
the representation or its training calibration becomes a specific suspect.
If both implementations fail under the same settings on a new target,
the result establishes that transfer failure, rather than a theorem against
the method. Calibration targets and untouched randomized targets must remain
separate throughout this comparison.

Finally, the sampling corrections retained in the literature explain how an
imperfect density can still be useful. Metropolis correction preserves the
target for an appropriate fixed proposal; importance and SMC corrections
recover target integrals under their support and integrability assumptions.
Their finite accuracy and efficiency still need evidence. Accurate corrected
sampling, accurate raw-flow density and Gaussian latent scores therefore
remain distinct conclusions. For the state-space application, the decisive
test is the frozen map's actual posterior computation, accompanied by the
distribution and geometric checks needed to explain its behavior.
